\chapter{Estado actual y próximos pasos}

\section{Estado funcional}
En el punto actual del proyecto ya se dispone de:
\begin{itemize}
  \item infraestructura local desplegada;
  \item conversación mediante Open WebUI;
  \item inferencia local con Ollama;
  \item memoria documental con Nextcloud, ingesta y Qdrant;
  \item control domótico mediante Home Assistant;
  \item búsqueda de fotos con Immich;
  \item consulta de cámaras con Frigate;
  \item notificaciones y monitorización mediante Telegram, Netdata y exportación a Home Assistant;
  \item repositorio privado con configuraciones saneadas y herramientas versionadas.
\end{itemize}

\section{Trabajo pendiente}
Antes de cerrar el TFG quedan varios frentes:
\begin{enumerate}
  \item completar capturas y evidencias visuales;
  \item ejecutar y verificar la política de backup con NVMe externo;
  \item terminar el dashboard único en Home Assistant;
  \item incorporar, si se dispone del medidor, datos reales de consumo y costes;
  \item consolidar pruebas de casos de uso completos;
  \item revisar la memoria final para que toda decisión relevante quede respaldada por fuentes y evidencias.
\end{enumerate}

\section{Valor del seguimiento}
El seguimiento confirma que el proyecto ha superado la fase de viabilidad y se encuentra en fase de consolidación. Los riesgos principales ya no son si Odín puede existir, sino cómo presentarlo con claridad, medirlo mejor y cerrar de forma responsable lo que todavía está abierto.
